\subsection*{沿着数据走一遍整体结构}
一段文字先分成 token，每个 token 查表得到嵌入向量，再加入或结合位置信息。只有内容、没有顺序，就无法区分同一组词的不同排列。随后经过多层注意力和前馈计算，向量逐渐携带上下文信息，最后投影成词表上每个候选的分数。

分数不是概率：它可以为负，也不必相加为一。Softmax 先指数化使其为正，再用总和归一化。相对分数差决定概率偏好，同加一个常数只引入共同因子，除法会将其抵消。

\subsection*{从一个两分类例子看出数值实现的意义}
分数 $(0,\ln3)$ 指数化后是 $(1,3)$，归一化得到 $(1/4,3/4)$。两分数都减去 $\ln3$ 后，虽然数值改变，概率不变。实现中常减去最大分数，避免指数过大；这个技巧基于精确恒等式，浮点误差仍需单独考虑。

嵌入空间的方向可能与某些语义变化相关，但不要把每个坐标轴都当成固定概念标签。向量表示允许特征分布在许多坐标中，网络也可能在不同层使用不同表示。

本节只建立全景路线：嵌入负责表示，注意力交换位置间信息，前馈层处理各位置的特征，输出投影把表示转为预测分数。真实架构还涉及归一化、残差等细节；理解它们时仍应沿数据形状和作用追踪，而不是只记一张盒子图。
